% ------------------------------------------------------------------------------
% Slide 7: El Problema del Punto de Partida: Reposo, Vocales Intermedias y Coarticulación
% ------------------------------------------------------------------------------
\begin{frame}{El Problema del Punto de Partida: Reposo, Vocales Intermedias y Coarticulación}
    \begin{columns}[c]
        \begin{column}{0.51\textwidth}
            \scriptsize
            \begin{block}{Dinámica Muscular y Coarticulación}
                El electromiograma registra la fuerza para ir de una postura a otra: $\Delta \vec{x}(t) = \vec{x}_{\text{final}} - \vec{x}_{\text{inicial}}$.
                \begin{itemize}\setlength{\itemsep}{1pt}\setlength{\parskip}{0pt}
                    \item Si el hablante no afloja la cara entre letras, la posición de partida cambia continuamente.
                    \item \textbf{Paso de /e/ a /i/:} El movimiento restante es muy chico frente al trayecto desde el reposo.
                    \item \textbf{Efecto en el modelo:} Una /i/ que parte desde la /e/ produce la misma señal que una /e/ desde el reposo.
                \end{itemize}
            \end{block}
            \vspace{0.1cm}
            \begin{exampleblock}{Estandarizar la Tarea en el Laboratorio}
                Definir una regla fija para los ensayos: exigir volver a relajar la cara antes de cada vocal, o alimentar al modelo con la historia de la vocal previa.
            \end{exampleblock}
        \end{column}
        \begin{column}{0.47\textwidth}
            \centering
            \placeholderfigura{0.70\textheight}{Esquema o traza temporal mostrando la diferencia de amplitud al ir de Reposo -> I vs. ir de E -> I sin volver al reposo}
        \end{column}
    \end{columns}
\end{frame}
